% TOP_PROBLEM: {"id": 41, "title": "Singular Cardinals Hypothesis", "category_id": 10, "category": {"id": 10, "name": "set_theory", "display_name": "Set Theory", "description": "Foundations of mathematics, infinite sets, and cardinality.", "slug": "set-theory", "order_index": 10, "created_at": "2026-07-31T15:26:25.671Z"}, "status": "open"}
% FIRST_POSTED: null
% ENTRY_KIND: statement_only
% Imported from: batch01.tex; UnsolvedMath ID 41
\documentclass[11pt]{article}
\usepackage[margin=25mm]{geometry}
\usepackage{fontspec}
\setmainfont{Latin Modern Roman}
\usepackage{amsmath,amssymb,mathtools}
\usepackage{unicode-math}
\setmathfont{Latin Modern Math}
\usepackage{xurl,hyperref}
\hypersetup{colorlinks=true,urlcolor=blue}
\setcounter{secnumdepth}{0}
\setlength{\parindent}{0pt}
\setlength{\parskip}{5pt}
\setlength{\emergencystretch}{3em}
\newcommand{\sref}[2]{\hyperlink{source:#1:#2}{[S#2]}}
\newcommand{\cataloguescope}[1]{\paragraph{Recorded target.}#1}
\begin{document}
% UnsolvedMath ID 41; source catalogue code SET-002
\setcounter{section}{40}
\section{Singular Cardinals Hypothesis}\label{41}
{\small\sffamily UnsolvedMath ID 41 · Set Theory}
\subsection{Definitions and mathematical statement}

\paragraph{Shared notation.}$\mathbb N=\{1,2,\ldots\}$, and $\mathbb Z,\mathbb Q,\mathbb R,\mathbb C$ denote the integers, rationals, reals and complex numbers. Any other conventions are specified locally.


Regard a cardinal $\kappa$ as its initial ordinal. Its cofinality $\operatorname{cf}(\kappa)$ is the least cardinality of a subset $A\subseteq\kappa$ unbounded in $\kappa$: for every $\alpha<\kappa$ there is $\beta\in A$ with $\alpha\le\beta$. An infinite cardinal is singular when $\operatorname{cf}(\kappa)<\kappa$. It is a strong limit when $2^\lambda<\kappa$ for every cardinal $\lambda<\kappa$, where $2^\lambda=|\mathcal P(\lambda)|$. Write $\kappa^+$ for the least cardinal greater than $\kappa$. The underlying axioms are ZFC.
\paragraph{Statement recorded in the supplied source.}
The singular cardinals hypothesis (SCH) asserts
\[
 \forall\kappa\quad
 [\kappa\text{ is an infinite singular strong limit cardinal}]
 \ \Longrightarrow\ 2^\kappa=\kappa^+.
\]
Is this universal cardinal-arithmetic assertion valid? \sref{41}{2}
The generalized continuum hypothesis, which asserts $2^\lambda=\lambda^+$ for every infinite cardinal $\lambda$, implies SCH. There are also relative-consistency constructions of failures of SCH from suitable large-cardinal hypotheses; those hypotheses are stronger than merely the consistency of ZFC. Source~S2 discusses both the assertion and the consistency strength of its failure. \sref{41}{2}
\subsection{Short English statement}
Does the power set of every singular strong limit cardinal have the smallest cardinality it could have, namely the next larger cardinal? Its behavior is sensitive to additional set-theoretic assumptions.
\subsection{Sources}
\begin{description}
\item[\hypertarget{source:41:1}{S1}] ULAM.AI and the UnsolvedMath Contributors, UnsolvedMath: A Curated Collection of Open Mathematics Problems, record 41 (SET-002). CC BY 4.0. \url{https://huggingface.co/datasets/ulamai/UnsolvedMath}.
\item[\hypertarget{source:41:2}{S2}] William J. Mitchell, On the Singular Cardinal Hypothesis, arXiv:math/9204202 (1992), Section 1, opening definition and Theorem 1.1. \url{https://arxiv.org/pdf/math/9204202}.
\end{description}
\label{end:41}
\end{document}
